% GENERATED FILE. Edit canonical lessons, then run npm run generate:books.
% canonical-source-hash: fnv1a64:d1b13ce697ea2c99
% canonical-lessons: TE-C413-sastram, TE-C413-sastravetta, TE-C413-prayogam, TE-C413-prayogasala, TE-C413-aviskarincu, TE-R413-review-406-413

\chapter{Science}
\label{ch:te-science}
\hlchaptermodality{413}

\begin{chapteropening}
\textbf{By the end of this chapter:} \emph{I can talk about science, scientists and experiments.}

Complete the last lesson of chapter 413: I can talk about science, scientists and experiments.
\end{chapteropening}

\section[\texorpdfstring{\te{శాస్త్రం} (śāstraṁ)}{శాస్త్రం (śāstraṁ)}]{\te{శాస్త్రం} (śāstraṁ) — a branch of learning; a science}

\label{lesson:TE-C413-sastram}



\begin{quote}
Before the new one: say the Telugu for information, then the Telugu for to explore.
\end{quote}

\subsection*{You'll want to know: \te{శాస్త్రం}}
\textbf{\te{శాస్త్రం}} — \emph{śāstraṁ} — \textquotedblleft{}a branch of learning; a science\textquotedblright{}.

\begin{quote}
\te{నాకు} \te{చిన్నప్పటి} \te{నుంచి} \te{శాస్త్రం} \te{అంటే} \te{ఇష్టం}.
\end{quote}

\begin{quote}
\emph{nāku cinnappaṭi nuñci śāstraṁ aṇṭē iṣṭaṁ.}
\end{quote}

\begin{quote}
I have liked science since I was small.
\end{quote}

Plural: \te{శాస్త్రాలు}, \emph{śāstrālu}.

In the names of subjects: \te{భౌతికశాస్త్రం}, physics; \te{జీవశాస్త్రం}, biology.

\begin{cousinweb}
\textbf{\te{శాస్త్రం}} (\emph{śāstraṁ}) is borrowed from Sanskrit \emph{śāstra}, \textquotedblleft{}a command\textquotedblright{}, and also \textquotedblleft{}advice, a treatise\textquotedblright{} (CDIAL 12422); Telugu uses it for a branch of knowledge.
\end{cousinweb}

\begin{grammarlens}[title={What you've built}]
One more word for giving reasons and opinions.
\end{grammarlens}

\subsection*{Your turn}
\begin{itemize}
\raggedright
  \item \emph{Say it:} \emph{śāstraṁ}
  \item \emph{Say it:} \emph{śāstraṁ}, once more
  \item \emph{Say it:} \emph{nāku cinnappaṭi nuñci śāstraṁ aṇṭē iṣṭaṁ.} — I have liked science since I was small.
  \item \emph{Recall:} say the Telugu for information, then the Telugu for to explore, then say \emph{śāstraṁ} again
\end{itemize}

\begin{tcolorbox}[breakable,colback=teal!4,colframe=teal!35!black,title={Before you move on}]
What does \te{శాస్త్రం} mean? (\textquotedblleft{}A branch of learning; a science\textquotedblright{}.) Say it once more.
\end{tcolorbox}
\section[\texorpdfstring{\te{శాస్త్రవేత్త} (śāstravētta)}{శాస్త్రవేత్త (śāstravētta)}]{\te{శాస్త్రవేత్త} (śāstravētta) — a scientist}

\label{lesson:TE-C413-sastravetta}



\begin{quote}
Before the new one: say the Telugu for to explore, then the Telugu for a branch of learning.
\end{quote}

\subsection*{You'll want to know: \te{శాస్త్రవేత్త}}
\textbf{\te{శాస్త్రవేత్త}} — \emph{śāstravētta} — \textquotedblleft{}a scientist\textquotedblright{}.

\begin{quote}
\te{ఆమె} \te{శాస్త్రవేత్త} \te{కావాలని} \te{అనుకుంటోంది}.
\end{quote}

\begin{quote}
\emph{āme śāstravētta kāvālani anukuṇṭōndi.}
\end{quote}

\begin{quote}
She wants to become a scientist.
\end{quote}

Plural: \te{శాస్త్రవేత్తలు}, \emph{śāstravēttalu}.

\te{శాస్త్ర}, science, + \te{వేత్త}, one who knows.

\begin{cousinweb}
\textbf{\te{శాస్త్రవేత్త}} (\emph{śāstravētta}) is built from Sanskrit \emph{śāstra}, \textquotedblleft{}a treatise, a body of knowledge\textquotedblright{} (CDIAL 12422), and \emph{vettṛ}, \textquotedblleft{}one who knows\textquotedblright{} (Monier-Williams).
\end{cousinweb}

\begin{grammarlens}[title={What you've built}]
One more word for giving reasons and opinions.
\end{grammarlens}

\subsection*{Your turn}
\begin{itemize}
\raggedright
  \item \emph{Say it:} \emph{śāstravētta}
  \item \emph{Say it:} \emph{śāstravētta}, once more
  \item \emph{Say it:} \emph{āme śāstravētta kāvālani anukuṇṭōndi.} — She wants to become a scientist.
  \item \emph{Recall:} say the Telugu for to explore, then the Telugu for a branch of learning, then say \emph{śāstravētta} again
\end{itemize}

\begin{tcolorbox}[breakable,colback=teal!4,colframe=teal!35!black,title={Before you move on}]
What does \te{శాస్త్రవేత్త} mean? (\textquotedblleft{}A scientist\textquotedblright{}.) Say it once more.
\end{tcolorbox}
\section[\texorpdfstring{\te{ప్రయోగం} (prayōgaṁ)}{ప్రయోగం (prayōgaṁ)}]{\te{ప్రయోగం} (prayōgaṁ) — an experiment}

\label{lesson:TE-C413-prayogam}



\begin{quote}
Before the new one: say the Telugu for a branch of learning, then the Telugu for a scientist.
\end{quote}

\subsection*{You'll want to know: \te{ప్రయోగం}}
\textbf{\te{ప్రయోగం}} — \emph{prayōgaṁ} — \textquotedblleft{}an experiment\textquotedblright{}.

\begin{quote}
\te{మేము} \te{తరగతిలో} \te{ఒక} \te{చిన్న} \te{ప్రయోగం} \te{చేశాము}.
\end{quote}

\begin{quote}
\emph{mēmu taragatilō oka cinna prayōgaṁ cēśāmu.}
\end{quote}

\begin{quote}
We did a small experiment in class.
\end{quote}

Plural: \te{ప్రయోగాలు}, \emph{prayōgālu}.

\te{ప్రయోగం} \te{చేయు}, to carry out an experiment.

\begin{cousinweb}
\textbf{\te{ప్రయోగం}} (\emph{prayōgaṁ}) is borrowed from Sanskrit \emph{prayoga}, \textquotedblleft{}joining together, connection\textquotedblright{} (Monier-Williams), from \emph{pra-} and \emph{yuj}, \textquotedblleft{}to join\textquotedblright{}; Telugu uses it for an experiment or a trial.
\end{cousinweb}

\begin{grammarlens}[title={What you've built}]
One more word for giving reasons and opinions.
\end{grammarlens}

\subsection*{Your turn}
\begin{itemize}
\raggedright
  \item \emph{Say it:} \emph{prayōgaṁ}
  \item \emph{Say it:} \emph{prayōgaṁ}, once more
  \item \emph{Say it:} \emph{mēmu taragatilō oka cinna prayōgaṁ cēśāmu.} — We did a small experiment in class.
  \item \emph{Recall:} say the Telugu for a branch of learning, then the Telugu for a scientist, then say \emph{prayōgaṁ} again
\end{itemize}

\begin{tcolorbox}[breakable,colback=teal!4,colframe=teal!35!black,title={Before you move on}]
What does \te{ప్రయోగం} mean? (\textquotedblleft{}An experiment\textquotedblright{}.) Say it once more.
\end{tcolorbox}
\section[\texorpdfstring{\te{ప్రయోగశాల} (prayōgaśāla)}{ప్రయోగశాల (prayōgaśāla)}]{\te{ప్రయోగశాల} (prayōgaśāla) — a laboratory}

\label{lesson:TE-C413-prayogasala}



\begin{quote}
Before the new one: say the Telugu for a scientist, then the Telugu for an experiment.
\end{quote}

\subsection*{You'll want to know: \te{ప్రయోగశాల}}
\textbf{\te{ప్రయోగశాల}} — \emph{prayōgaśāla} — \textquotedblleft{}a laboratory\textquotedblright{}.

\begin{quote}
\te{మా} \te{కళాశాలలో} \te{కొత్త} \te{ప్రయోగశాల} \te{కట్టారు}.
\end{quote}

\begin{quote}
\emph{mā kaḷāśālalō kotta prayōgaśāla kaṭṭāru.}
\end{quote}

\begin{quote}
They built a new laboratory at our college.
\end{quote}

\te{ప్రయోగ}, experiment, + \te{శాల}, a hall, as in \te{కళాశాల}.

\begin{cousinweb}
\textbf{\te{ప్రయోగశాల}} (\emph{prayōgaśāla}) is built from Sanskrit \emph{prayoga} (Monier-Williams) and \emph{śālā}, \textquotedblleft{}a shed, a house, a hall\textquotedblright{} (CDIAL 12414), the ending in \te{కళాశాల}.
\end{cousinweb}

\begin{grammarlens}[title={What you've built}]
One more word for giving reasons and opinions.
\end{grammarlens}

\subsection*{Your turn}
\begin{itemize}
\raggedright
  \item \emph{Say it:} \emph{prayōgaśāla}
  \item \emph{Say it:} \emph{prayōgaśāla}, once more
  \item \emph{Say it:} \emph{mā kaḷāśālalō kotta prayōgaśāla kaṭṭāru.} — They built a new laboratory at our college.
  \item \emph{Recall:} say the Telugu for a scientist, then the Telugu for an experiment, then say \emph{prayōgaśāla} again
\end{itemize}

\begin{tcolorbox}[breakable,colback=teal!4,colframe=teal!35!black,title={Before you move on}]
What does \te{ప్రయోగశాల} mean? (\textquotedblleft{}A laboratory\textquotedblright{}.) Say it once more.
\end{tcolorbox}
\section[\texorpdfstring{\te{ఆవిష్కరించు} (āviṣkariñcu)}{ఆవిష్కరించు (āviṣkariñcu)}]{\te{ఆవిష్కరించు} (āviṣkariñcu) — to invent; to unveil}

\label{lesson:TE-C413-aviskarincu}



\begin{quote}
Before the new one: say the Telugu for an experiment, then the Telugu for a laboratory.
\end{quote}

\subsection*{You'll want to know: \te{ఆవిష్కరించు}}
\textbf{\te{ఆవిష్కరించు}} — \emph{āviṣkariñcu} — \textquotedblleft{}to invent; to unveil\textquotedblright{}.

\begin{quote}
\te{ఆ} \te{శాస్త్రవేత్త} \te{ఒక} \te{కొత్త} \te{మందును} \te{ఆవిష్కరించాడు}.
\end{quote}

\begin{quote}
\emph{ā śāstravētta oka kotta mandunu āviṣkariñcāḍu.}
\end{quote}

\begin{quote}
That scientist invented a new medicine.
\end{quote}

Past (I …): \te{ఆవిష్కరించాను}, \emph{āviṣkariñcānu}.

Also to unveil a statue or launch a book: \te{పుస్తకాన్ని} \te{ఆవిష్కరించు}.

\begin{cousinweb}
\textbf{\te{ఆవిష్కరించు}} (\emph{āviṣkariñcu}) is built on Sanskrit \emph{āviṣkaraṇa}, \textquotedblleft{}making visible, manifestation\textquotedblright{} (Monier-Williams), from \emph{āviṣkṛ}, \textquotedblleft{}to make apparent, to reveal\textquotedblright{}, with the Telugu verb ending -\te{ఇంచు}.
\end{cousinweb}

\begin{grammarlens}[title={What you've built}]
One more word for giving reasons and opinions.
\end{grammarlens}

\subsection*{Your turn}
\begin{itemize}
\raggedright
  \item \emph{Say it:} \emph{āviṣkariñcu}
  \item \emph{Say it:} \emph{āviṣkariñcu}, once more
  \item \emph{Say it:} \emph{ā śāstravētta oka kotta mandunu āviṣkariñcāḍu.} — That scientist invented a new medicine.
  \item \emph{Recall:} say the Telugu for an experiment, then the Telugu for a laboratory, then say \emph{āviṣkariñcu} again
\end{itemize}

\begin{tcolorbox}[breakable,colback=teal!4,colframe=teal!35!black,title={Before you move on}]
What does \te{ఆవిష్కరించు} mean? (\textquotedblleft{}To invent; to unveil\textquotedblright{}.) Say it once more.
\end{tcolorbox}
\section[(dialogue)]{Words from chapters 406–413 — a review}

\label{lesson:TE-R413-review-406-413}



\begin{quote}
Say the Telugu for corruption and a bribe.
\end{quote}

\subsection*{Your turn}
Say these aloud:

\begin{itemize}
\raggedright
  \item corruption, a bribe, transparency, an allegation, to allege
  \item a protest, a movement, to condemn, a claim, a reform
  \item society, social, an association, a member, discrimination
  \item violence, harmony, unity, solidarity, to put into practice
  \item service, a donation, voluntary, an activist, a donor
  \item a census, a survey, statistics, public opinion, a percentage
  \item research, to research, a researcher, information, to explore
  \item a branch of learning, a scientist, an experiment, a laboratory, to invent
\end{itemize}

\begin{tcolorbox}[breakable,colback=teal!4,colframe=teal!35!black,title={Before you move on}]
Corruption? (\textbf{\te{అవినీతి}}, \emph{avinīti}.) A bribe? (\textbf{\te{లంచం}}, \emph{lañcaṁ}.) Transparency? (\textbf{\te{పారదర్శకత}}, \emph{pāradarśakata}.) An allegation? (\textbf{\te{ఆరోపణ}}, \emph{ārōpaṇa}.) To allege? (\textbf{\te{ఆరోపించు}}, \emph{ārōpiñcu}.) A protest? (\textbf{\te{నిరసన}}, \emph{nirasana}.) A movement? (\textbf{\te{ఉద్యమం}}, \emph{udyamaṁ}.) To condemn? (\textbf{\te{ఖండించు}}, \emph{khaṇḍiñcu}.) A claim? (\textbf{\te{డిమాండు}}, \emph{ḍimāṇḍu}.) A reform? (\textbf{\te{సంస్కరణ}}, \emph{saṁskaraṇa}.) Society? (\textbf{\te{సమాజం}}, \emph{samājaṁ}.) Social? (\textbf{\te{సామాజిక}}, \emph{sāmājika}.) An association? (\textbf{\te{సంఘం}}, \emph{saṅghaṁ}.) A member? (\textbf{\te{సభ్యుడు}}, \emph{sabhyuḍu}.) Discrimination? (\textbf{\te{వివక్ష}}, \emph{vivakṣa}.) Violence? (\textbf{\te{హింస}}, \emph{hiṁsa}.) Harmony? (\textbf{\te{సామరస్యం}}, \emph{sāmarasyaṁ}.) Unity? (\textbf{\te{ఐక్యత}}, \emph{aikyata}.) Solidarity? (\textbf{\te{సంఘీభావం}}, \emph{saṅghībhāvaṁ}.) To put into practice? (\textbf{\te{ఆచరించు}}, \emph{ācariñcu}.) Service? (\textbf{\te{సేవ}}, \emph{sēva}.) A donation? (\textbf{\te{విరాళం}}, \emph{virāḷaṁ}.) Voluntary? (\textbf{\te{స్వచ్ఛంద}}, \emph{svacchanda}.) An activist? (\textbf{\te{కార్యకర్త}}, \emph{kāryakarta}.) A donor? (\textbf{\te{దాత}}, \emph{dāta}.) A census? (\textbf{\te{జనగణన}}, \emph{janagaṇana}.) A survey? (\textbf{\te{సర్వే}}, \emph{sarvē}.) Statistics? (\textbf{\te{గణాంకాలు}}, \emph{gaṇāṅkālu}.) Public opinion? (\textbf{\te{ప్రజాభిప్రాయం}}, \emph{prajābhiprāyaṁ}.) A percentage? (\textbf{\te{శాతం}}, \emph{śātaṁ}.) Research? (\textbf{\te{పరిశోధన}}, \emph{pariśōdhana}.) To research? (\textbf{\te{పరిశోధించు}}, \emph{pariśōdhiñcu}.) A researcher? (\textbf{\te{పరిశోధకుడు}}, \emph{pariśōdhakuḍu}.) Information? (\textbf{\te{సమాచారం}}, \emph{samācāraṁ}.) To explore? (\textbf{\te{అన్వేషించు}}, \emph{anvēṣiñcu}.) A branch of learning? (\textbf{\te{శాస్త్రం}}, \emph{śāstraṁ}.) A scientist? (\textbf{\te{శాస్త్రవేత్త}}, \emph{śāstravētta}.) An experiment? (\textbf{\te{ప్రయోగం}}, \emph{prayōgaṁ}.) A laboratory? (\textbf{\te{ప్రయోగశాల}}, \emph{prayōgaśāla}.) To invent? (\textbf{\te{ఆవిష్కరించు}}, \emph{āviṣkariñcu}.)
\end{tcolorbox}
